% Insert once before the odd-root catalogue.  The path-calculus inputs are
% q3:lem:spectral and q3:app:foundations; no new kernel hypothesis is used.
\subsection{Finite calculus for the odd-root estimates}
\label{sec:odd-common-calculus}

Work on a finite probability space after deleting zero-mass atoms. Let
\(B=-K\), where \(K\) is a real antisymmetric kernel with \(|K|\le1\),
and write
\[
 T=J+B,\quad A=J-B=T^*,\quad u=B\one,\quad p=1-u,\quad q=1+u,
 \qquad d=\E u^2,\quad \delta_j=\|B^j\one\|_2^2.
\]
The four arm messages are
\(f_n=A^n\one\), \(\bar f_n=T^n\one\),
\(h_n=Tf_{n-1}\) and \(W_n=A\bar f_{n-1}\).
Thus \(f_n\) points outward, \(\bar f_n\) inward, \(h_n\) first
inward and then outward, and \(W_n\) first outward and then inward.
The path calculus in Section~\ref{q3:app:foundations} gives
\begin{equation}\label{eq:odd-common:inputs}
\begin{gathered}
 0\le d\le\tfrac13,\quad d^2\le\delta_2\le\theta d,
 \quad\delta_{j+1}\le\theta\delta_j,
 \quad\|B\|^2\le\theta=\tfrac{13}{32},\quad\kappa=\tfrac{19}{32},\\
 |BV(x)|\le\|V\|_2,\qquad
 0\le f_n,\bar f_n,h_n,W_n\le2,\qquad
 \E f_n=\E\bar f_n\le1,\\
 0\le f_nh_n,\bar f_nW_n\le1.
\end{gathered}
\end{equation}
Here the path bounds concern \(n\ge1\), and the mean bound also holds
at \(n=0\). The last line permits only genuine equal-length pair
decorations. All operators and expectations below use the original
measure. Since \(p(x),q(x)\ge\mu(x)>0\), actual tree messages and the
upper bounds dominating them are positive. When \(d=0\), leaf
integration gives density one; subsequent estimates take \(d>0\).

\begin{lemma}[Finite spectral coordinates]\label{lem:odd-common:spectral}
Let \(G_n(z)=\sum_{j=0}^{n-1}z^j\), \(G_0=0\), and
\[
 \beta_n=\sum_{j=1}^{\lfloor(n-1)/2\rfloor}(-1)^{j-1}\delta_j
 \quad(n\ge1).
\]
Empty sums are zero. With \(\tau=d-\delta_2\),
\begin{equation}\label{eq:Q188:conv}
 S_n:=\one+G_n(B)u
   =\bar f_n+\sum_{j=3}^{n}\beta_j\bar f_{n-j},\qquad
 \tau\le\beta_j\le d\quad(j\ge3),\qquad
 \kappa d\le\tau\le\tfrac29.
\end{equation}
There is a positive measure \(\nu\) on \([0,\theta]\) of mass \(d\)
such that \(\int x^j\,d\nu=\delta_{j+1}\). If
\(P(i\sqrt x)=a(x)+i\sqrt x\,b(x)\), then
\begin{equation}\label{eq:odd-common:coordinates}
\begin{aligned}
 \E[P(B)u]&=-\int b\,d\nu,&
 \E[BP(B)u]&=-\int a\,d\nu,\\
 \|P(B)u+c\one\|_2^2
   &=\int(a^2+xb^2)\,d\nu-2c\int b\,d\nu+c^2.
\end{aligned}
\end{equation}
For \(G_m=E+zC\), \((E,C)=(1,0)\) at \(m=1\). For \(m\ge2\),
\(1-x\le E,C\le1\); even \(m\) has \(E=C\), and odd \(m\ge3\)
has \(E=1-xC\). In particular \(G_2=1+z\) and
\(|G_2|^2=1+x\); no contraction bound is imposed at this index.
For \(m\ge3\), \(E\le1-x+x^2\) and \(|G_m|\le1\).
\end{lemma}

\begin{proof}
Skew-adjointness gives
\(\E B^{2j+1}\one=0\) and
\(\E B^{2j}\one=(-1)^j\delta_j\).
Hence \(TS_n=S_{n+1}-\beta_{n+1}\one\), which proves the finite
convolution by induction. The alternating sums lie between
\(d-\delta_2\) and \(d\). Also
\(\tau\le d-d^2\le2/9\).
The spectral measure of \(-B^2\) at \(u\) gives \(\nu\) and the norm
identity. Skew-adjointness then gives the two means and the displayed
constant cross term. The coordinate relations follow directly by
separating the finite even and odd powers of \(G_m\).
For odd \(m\ge3\),
\(|G_m|^2=(1-xC)^2+xC^2\le1\).
For even \(m\ge4\), use \(C\le1-x+x^2\) and
\[
 1-(1+x)(1-x+x^2)^2
 =x(1-x-x^2+x^3-x^4)\ge3x/16\ge0
 \quad(0\le x\le1/2).
\]
\end{proof}

\begin{lemma}[One finite mask]\label{lem:odd-common:mask}
For every finite polynomial \(P(z)=\sum_j\epsilon_jz^j\),
\(\epsilon_j\in\{0,1\}\), and every \(0\le x\le\theta\),
\begin{equation}\label{eq:odd-common:mask}
 |P(i\sqrt x)|^2\le\frac{1+x}{(1-x^2)^2}
 \le\frac{32768}{16245}<\frac{101}{50}.
\end{equation}
Consequently, for every integer \(r\ge0\),
\(\|B^rP(\pm B)u\|^2\le(101/50)\theta^r d\).
The constants are independent of the number of nonzero coefficients.
\end{lemma}
\begin{proof}
In each real or imaginary coordinate, separate the two alternating
signs and take the larger same-sign subseries. Each such series is
bounded by \((1-x^2)^{-1}\). This proves the first inequality;
its two positive factors increase on \([0,\theta]\).
Integrate against the original spectral measure for the norm bound.
The geometric series bound finite sums only; no path or message is
replaced by an infinite one.
\end{proof}

\begin{lemma}[A joint finite path bound]\label{lem:odd-common:joint-paths}
For every \(n\ge3\), \(m\ge1\) and \(0\le x\le\theta\),
\begin{equation}\label{eq:odd-common:joint-paths}
 |G_n(z)-zG_m(z)|\le\tfrac65,\qquad z=i\sqrt x.
\end{equation}
When \(m=1\), the bound improves to one.
\end{lemma}
\begin{proof}
The difference is \(1+z^{m+1}G_{n-1-m}\) if \(n-1\ge m\), and
\(1-z^nG_{m-n+1}\) otherwise. The empty tail gives one.
For a nonempty tail, the finite coordinate bounds above give the
following bounds for the square:
\[
\begin{array}{c|c}
 \text{tail term}&\text{squared upper bound}\\ \hline
 +z^2G_l&(1-x+x^2)^2+x^3\le1\\
 -z^3G_l&1+x^3\\
 +z^3G_l&1+2x^2+x^3+x^4\\
 \pm z^kG_l,\ k\ge4&1+2x^2+x^4(1+x).
\end{array}
\]
The third row is at most \(1493553/1048576<36/25\) at
\(x=\theta\); the last is at most
\(345857/248832<7/5<36/25\) even at \(x=5/12\).
The second row is smaller still. For the first row,
the difference from one is
\(x[-2+3x-x^2+x^3]\le0\) on \([0,1/2]\).
When \(m=1\), use this first row on
\(G_n-z=1+z^2G_{n-2}\). These bounds include \(G_1\) and \(G_2\).
\end{proof}
